\section{Fiber checkpoints}

\label{sec:checkpoints}
%% ============================================================================

\subsection{Why incremental repair}

Repairs in 3.0~\cite{lp010-quasarstm-3} re-execute the entire
transaction from \texttt{pc=0}. For long-running transactions---hot
routers, multi-hop swaps, batched DEX matches---this is wasteful: the
conflicting read often happens deep into the transaction, after many
side-effect-free arithmetic and stack operations. 4.0 ships
checkpoint-based incremental repair as the production path.

\subsection{Checkpoint record}

\begin{lstlisting}[language=C++]
struct FiberCheckpoint {
    uint32_t fiber_id;             // index in the fiber pool
    uint32_t pc;                   // EVM program counter at checkpoint
    uint64_t gas_remaining;
    uint16_t stack_depth;
    uint16_t mem_size;
    uint8_t  call_depth;           // CALL frames opened at checkpoint
    uint8_t  pad[3];
    Hash     state_digest;         // commitment over MVCC reads since start
    uint32_t rwset_offset;         // index into rwset_arena
    uint16_t rwset_count;
    uint16_t pad2;
};
\end{lstlisting}

A checkpoint is emitted at every \texttt{CALL}-frame return, every
explicit checkpoint opcode (\texttt{CKPT} pseudo-opcode for
precompiles), and every cold-state suspend (\S\S 7.4 of
3.0~\cite{lp010-quasarstm-3}). Each fiber maintains a ring of the last
$N$ checkpoints (default $N=4$); older checkpoints are recycled into
the per-round arena.

\subsection{Repair from checkpoint}

When Tier~1 / Tier~2 validation fails for a transaction that has emitted
checkpoints, repair restores the latest checkpoint $c$ such that
$c.\mathit{pc}$ is strictly less than the program counter of the
conflicting read, then re-executes from $c.\mathit{pc}$. The fiber's
stack, memory, gas, and call-depth are restored from $c$; the read/write
set since $c$ is truncated.

\begin{algorithm}[ht]
\caption{\texttt{repair\_from\_checkpoint($t$)}}
\label{alg:repair-from-checkpoint}
\begin{algorithmic}[1]
\State $\mathrm{pc}^* \gets$ program counter of conflicting read in $t$
\State $c \gets$ latest checkpoint with $c.\mathit{pc} < \mathrm{pc}^*$
\If{no such $c$}
  \State \Return \texttt{repair\_from\_start}($t$) \Comment{3.0 path}
\EndIf
\State restore $\mathit{stack}$, $\mathit{mem}$, $\mathit{gas}$, $\mathit{call\_depth}$ from $c$
\State truncate $t.\mathit{rwset}$ to entries $\le c.\mathit{rwset\_offset}$
\State bump $t.\mathit{incarnation}$
\State enqueue $t$ at \texttt{ServiceId::Exec} with $\mathit{pc} \gets c.\mathit{pc}$
\end{algorithmic}
\end{algorithm}

\subsection{Cost reduction}

Measured on the cevm router workload (10-leg multi-hop swap, 1024
concurrent transactions, single hot pool):

\begin{itemize}[leftmargin=2em]
\item 3.0 (full re-execute): mean repair cost $\approx 240\,\mu\text{s}$ per repair.
\item 4.0 (checkpoint-based): mean repair cost $\approx 9$--$48\,\mu\text{s}$
      per repair, depending on which CALL-frame the conflict landed
      in. \textbf{5--30$\times$ reduction.}
\end{itemize}

\subsection{Storage cost}

Each checkpoint costs $\approx 96\text{B}$ of metadata plus a snapshot
of the live MVCC read set since the previous checkpoint. With the ring
size $N=4$ and a typical fiber accumulating $\le 8$ MVCC reads per
checkpoint window, total per-fiber checkpoint cost is $\approx 1\text{kB}$,
allocated from the per-round arena and reset on \texttt{end\_round()}.

\subsection{Determinism}

The checkpoint emit set is a function of the transaction's bytecode
and the canonical-order schedule (\texttt{CALL}-frame returns are
deterministic in the EVM trace). Two validators executing the same
canonical order emit the same checkpoint sequence per fiber, so the
\emph{repair-from-checkpoint} restoration is deterministic.

\begin{lemma}[Checkpoint repair preserves determinism]
\label{lem:checkpoint-det}
For any transaction $t$ that fails Tier~1 / Tier~2 validation in a
canonical-order trace, \texttt{repair\_from\_checkpoint}($t$) yields
the same final committed state as \texttt{repair\_from\_start}($t$).
\end{lemma}
\begin{proof}[Proof sketch]
The checkpointed state $(\mathit{stack}, \mathit{mem}, \mathit{gas},
\mathit{call\_depth})$ is the value $t$ had at $\mathit{pc}=c.\mathit{pc}$
during its first execution. Re-execution from $c.\mathit{pc}$ proceeds
through the same instruction sequence as the suffix
$[c.\mathit{pc}, \mathit{end}]$ of the original execution---except
that the conflicting read now sees the new visible version. The MVCC
visibility rule under canonical order (3.0 \S\,5) is the same; the
new visible version is canonical; hence the repaired result is the
same as a from-scratch re-execution that observes the same visible
versions for all reads.
\end{proof}

%% ============================================================================
